\chapter{Configuration and Mesh Independence of Validation Geometries}
\label{app:validation_configurations}

\providecommand{\meshplaceholder}[1]{%
	\fbox{\parbox[c][0.20\textheight][c]{0.88\linewidth}{\centering \textbf{Placeholder}\\[0.5em]#1}}%
}

To ensure complete scientific reproducibility and support the cross-code validation campaign detailed in Chapter~\ref{ch:benchmark_results}, this appendix documents the comprehensive physical boundary conditions, numerical FEniCS configurations, and Ansys spatial grid independence studies for all four topology optimization benchmarks. Because the native FEniCS topology optimization utilizes a fixed, Brinkman-penalized Cartesian grid, the transition to a sharp, body-fitted CAD geometry in Ansys introduces potential spatial discretization errors. Establishing grid independence ensures that the reported pressure drops ($\Delta p$) accurately reflect the underlying solver physics and turbulence models, rather than numerical remeshing noise.

\section{Near-Wall Treatment and Meshing Strategy}
\label{sec:app_mesh_strategy}
The extracted validation geometries must satisfy the strict near-wall discretization requirements of the selected commercial turbulence closures. Both the Spalart-Allmaras (SA) and SST $k$-$\omega$ models utilized in the Ansys Fluent re-analysis integrate the governing transport equations completely through the viscous sublayer to the solid boundary. This approach strictly circumvents the use of empirical wall functions, necessitating a highly controlled discretization of the boundary layer.

For all evaluated geometries across every benchmark, a hybrid meshing strategy was deployed in Ansys. The bulk flow domains were discretized utilizing unstructured polyhedral elements to efficiently manage complex, optimized internal curves. At all extracted solid boundaries ($\bar{\gamma} = 0$), explicit inflation layers were generated. To guarantee accurate resolution of the severe velocity gradients within the viscous sublayer, the first-layer cell height ($\Delta y_1$) was calibrated to maintain a non-dimensional wall distance of $y^+ \le 1$ across all solid surfaces. Subsequent prism layers were generated with a controlled geometric growth rate not exceeding $1.2$ to ensure a smooth volumetric transition into the bulk mesh.

% ---------------------------------------------------------
% BENCHMARK 1: DIFFUSER
% ---------------------------------------------------------
\section{Benchmark 1: Borrvall Diffuser}
\label{app:config_diffuser}

The diffuser benchmark evaluates the framework's ability to mitigate adverse pressure gradients in expanding flows. The geometric configuration and operational fluid parameters defining this design domain are detailed in Table~\ref{tab:app_diffuser_physical}.

\begin{table}[htbp]
	\centering
	\footnotesize
	\caption{Geometric, physical, and boundary-condition parameters for the Diffuser benchmark.}
	\label{tab:app_diffuser_physical}
	\renewcommand{\arraystretch}{1.2}
	\begin{tabularx}{\textwidth}{l X}
		\hline
		\textbf{Parameter} & \textbf{Diffuser setting} \\
		\hline
		Design Domain & $L \times L$ (with $L = 1.0$ m) \\
		Non-design fluid & Inlet extension: $0.2L$; Outlet extension: $0.2L$ \\
		Operational regime & Turbulent, $Re = 1{,}000$ \\
		Inlet condition & Uniform horizontal velocity, prescribed turbulence intensity \\
		Outlet condition & Relative static pressure $p = 0$ Pa \\
		Wall conditions & No-slip velocity ($\mathbf{u} = 0$), zero eddy viscosity ($\tilde{\nu} = 0$) \\
		\hline
	\end{tabularx}
\end{table}

To execute the optimization, specific penalty schedules and volume constraints were strictly enforced within the custom finite element solver. These numerical settings are summarized in Table~\ref{tab:app_diffuser_numerical}.

\begin{table}[htbp]
	\centering
	\footnotesize
	\caption{Numerical FEniCS solver settings and optimization parameters for the Diffuser.}
	\label{tab:app_diffuser_numerical}
	\renewcommand{\arraystretch}{1.2}
	\begin{tabularx}{\textwidth}{l c X}
		\hline
		\textbf{Parameter} & \textbf{Symbol} & \textbf{Setting} \\
		\hline
		Velocity space & $\mathbf{u}$ & Vector CG2 \\
		Pressure space & $p$ & Scalar CG1 \\
		SA working-variable space & $\tilde{\nu}$ & Scalar CG1 \\
		Adjoint formulation & -- & Frozen turbulence assumption \\
		Volume constraint & $V_{max}$ & [TBD]\% \\
		RAMP penalization factor & $q$ & [TBD] \\
		\hline
	\end{tabularx}
\end{table}

Upon convergence, the geometry was extracted and imported into Ansys Fluent. A formal grid independence study was conducted across three distinct spatial resolutions to verify convergence of the target metric. The results are recorded in Table~\ref{tab:app_diffuser_mesh}.

\begin{table}[htbp]
	\centering
	\footnotesize
	\caption{Grid convergence study for the extracted Diffuser geometry (Ansys Fluent).}
	\label{tab:app_diffuser_mesh}
	\renewcommand{\arraystretch}{1.2}
	\begin{tabularx}{\textwidth}{l c c c X}
		\hline
		\textbf{Model} & \textbf{Coarse Elements} & \textbf{Medium Elements} & \textbf{Fine Elements} & \textbf{Relative $\Delta p$ Change (Med-Fine)} \\
		\hline
		SA & [TBD] & [TBD] & [TBD] & [TBD]\% \\
		SST $k$-$\omega$ & [TBD] & [TBD] & [TBD] & [TBD]\% \\
		\hline
	\end{tabularx}
\end{table}

% ---------------------------------------------------------
% BENCHMARK 2: DOUBLE PIPE
% ---------------------------------------------------------
\section{Benchmark 2: Borrvall Double Pipe}
\label{app:config_doublepipe}

The double pipe benchmark forces complex fluid splitting and mixing. The domain setup, utilizing symmetric inlet and outlet boundary profiles, is formally outlined in Table~\ref{tab:app_doublepipe_physical}.

\begin{table}[htbp]
	\centering
	\footnotesize
	\caption{Geometric, physical, and boundary-condition parameters for the Double Pipe benchmark.}
	\label{tab:app_doublepipe_physical}
	\renewcommand{\arraystretch}{1.2}
	\begin{tabularx}{\textwidth}{l X}
		\hline
		\textbf{Parameter} & \textbf{Double Pipe setting} \\
		\hline
		Design Domain & $1.5L \times 1.0L$ (with $L = 1.0$ m) \\
		Non-design fluid & Inlet extensions: $0.2L$; Outlet extensions: $0.2L$ \\
		Operational regime & Turbulent, $Re = 1{,}660$ \\
		Inlet conditions & Two symmetric ports; uniform horizontal velocity \\
		Outlet conditions & Two symmetric ports; relative static pressure $p = 0$ Pa \\
		Wall conditions & No-slip velocity ($\mathbf{u} = 0$), zero eddy viscosity ($\tilde{\nu} = 0$) \\
		\hline
	\end{tabularx}
\end{table}

The optimization engine relied on the frozen turbulence adjoint subject to a restricted fluid volume fraction. The continuous Galerkin function spaces and penalty tuning parameters are provided in Table~\ref{tab:app_doublepipe_numerical}.

\begin{table}[htbp]
	\centering
	\footnotesize
	\caption{Numerical FEniCS solver settings and optimization parameters for the Double Pipe.}
	\label{tab:app_doublepipe_numerical}
	\renewcommand{\arraystretch}{1.2}
	\begin{tabularx}{\textwidth}{l c X}
		\hline
		\textbf{Parameter} & \textbf{Symbol} & \textbf{Setting} \\
		\hline
		Velocity space & $\mathbf{u}$ & Vector CG2 \\
		Pressure space & $p$ & Scalar CG1 \\
		Adjoint formulation & -- & Frozen turbulence assumption \\
		Volume constraint & $V_{max}$ & [TBD]\% \\
		RAMP penalization factor & $q$ & [TBD] \\
		\hline
	\end{tabularx}
\end{table}

Following extraction, the resulting CAD manifold was meshed with boundary layer inflation strategies. The macroscopic pressure drop convergence across successive mesh refinements is verified in Table~\ref{tab:app_doublepipe_mesh}.

\begin{table}[htbp]
	\centering
	\footnotesize
	\caption{Grid convergence study for the extracted Double Pipe geometry (Ansys Fluent).}
	\label{tab:app_doublepipe_mesh}
	\renewcommand{\arraystretch}{1.2}
	\begin{tabularx}{\textwidth}{l c c c X}
		\hline
		\textbf{Model} & \textbf{Coarse Elements} & \textbf{Medium Elements} & \textbf{Fine Elements} & \textbf{Relative $\Delta p$ Change (Med-Fine)} \\
		\hline
		SA & [TBD] & [TBD] & [TBD] & [TBD]\% \\
		SST $k$-$\omega$ & [TBD] & [TBD] & [TBD] & [TBD]\% \\
		\hline
	\end{tabularx}
\end{table}

% ---------------------------------------------------------
% BENCHMARK 3: PIPE BEND
% ---------------------------------------------------------
\section{Benchmark 3: Alexandersen Pipe Bend}
\label{app:config_pipebend}

The pipe bend isolates curvature-driven secondary flow phenomena. The specific geometry defining the $90^\circ$ deflection and the boundary conditions generating the high Dean numbers are presented in Table~\ref{tab:app_pipebend_physical}.

\begin{table}[htbp]
	\centering
	\footnotesize
	\caption{Geometric, physical, and boundary-condition parameters for the Pipe Bend benchmark.}
	\label{tab:app_pipebend_physical}
	\renewcommand{\arraystretch}{1.2}
	\begin{tabularx}{\textwidth}{l X}
		\hline
		\textbf{Parameter} & \textbf{Pipe Bend setting} \\
		\hline
		Design Domain & $90^\circ$ bend configuration \\
		Non-design fluid & Maintained at inlet and outlet to allow shear development \\
		Operational regime & Turbulent, High curvature ($De = \text{[TBD]}$) \\
		Inlet condition & Uniform velocity profile, prescribed turbulence intensity \\
		Outlet condition & Relative static pressure $p = 0$ Pa \\
		Wall conditions & No-slip velocity ($\mathbf{u} = 0$), zero eddy viscosity ($\tilde{\nu} = 0$) \\
		\hline
	\end{tabularx}
\end{table}

The numerical architecture utilized to solve this highly convective state within the FEniCS optimization loop is documented in Table~\ref{tab:app_pipebend_numerical}.

\begin{table}[htbp]
	\centering
	\footnotesize
	\caption{Numerical FEniCS solver settings and optimization parameters for the Pipe Bend.}
	\label{tab:app_pipebend_numerical}
	\renewcommand{\arraystretch}{1.2}
	\begin{tabularx}{\textwidth}{l c X}
		\hline
		\textbf{Parameter} & \textbf{Symbol} & \textbf{Setting} \\
		\hline
		Velocity space & $\mathbf{u}$ & Vector CG2 \\
		Pressure space & $p$ & Scalar CG1 \\
		Adjoint formulation & -- & Frozen turbulence assumption \\
		Volume constraint & $V_{max}$ & [TBD]\% \\
		\hline
	\end{tabularx}
\end{table}

The spatial discretization of the extracted curvature heavily impacts the Ansys solver physics. The asymptotic stabilization of the pressure drop across refined element counts is confirmed in Table~\ref{tab:app_pipebend_mesh}.

\begin{table}[htbp]
	\centering
	\footnotesize
	\caption{Grid convergence study for the extracted Pipe Bend geometry (Ansys Fluent).}
	\label{tab:app_pipebend_mesh}
	\renewcommand{\arraystretch}{1.2}
	\begin{tabularx}{\textwidth}{l c c c X}
		\hline
		\textbf{Model} & \textbf{Coarse Elements} & \textbf{Medium Elements} & \textbf{Fine Elements} & \textbf{Relative $\Delta p$ Change (Med-Fine)} \\
		\hline
		SA & [TBD] & [TBD] & [TBD] & [TBD]\% \\
		SST $k$-$\omega$ & [TBD] & [TBD] & [TBD] & [TBD]\% \\
		\hline
	\end{tabularx}
\end{table}

% ---------------------------------------------------------
% BENCHMARK 4: U-BEND
% ---------------------------------------------------------
\section{Benchmark 4: Alexandersen U-Bend}
\label{app:config_ubend}

The U-bend represents the most severe aerodynamic challenge, combining strong curvature with forced momentum reversal. The geometric restrictions placed on the solver to manage flow separation are outlined in Table~\ref{tab:app_ubend_physical}.

\begin{table}[htbp]
	\centering
	\footnotesize
	\caption{Geometric, physical, and boundary-condition parameters for the U-Bend benchmark.}
	\label{tab:app_ubend_physical}
	\renewcommand{\arraystretch}{1.2}
	\begin{tabularx}{\textwidth}{l X}
		\hline
		\textbf{Parameter} & \textbf{U-Bend setting} \\
		\hline
		Design Domain & $180^\circ$ flow reversal configuration \\
		Non-design fluid & Extended straight legs at inflow/outflow to prevent boundary intersection with separation zones \\
		Operational regime & Turbulent, massive separation \\
		Inlet condition & Uniform velocity profile, prescribed turbulence intensity \\
		Outlet condition & Relative static pressure $p = 0$ Pa \\
		Wall conditions & No-slip velocity ($\mathbf{u} = 0$), zero eddy viscosity ($\tilde{\nu} = 0$) \\
		\hline
	\end{tabularx}
\end{table}

The mathematical stabilization mechanisms and design constraints required to successfully optimize this separating flow in FEniCS are detailed in Table~\ref{tab:app_ubend_numerical}.

\begin{table}[htbp]
	\centering
	\footnotesize
	\caption{Numerical FEniCS solver settings and optimization parameters for the U-Bend.}
	\label{tab:app_ubend_numerical}
	\renewcommand{\arraystretch}{1.2}
	\begin{tabularx}{\textwidth}{l c X}
		\hline
		\textbf{Parameter} & \textbf{Symbol} & \textbf{Setting} \\
		\hline
		Velocity space & $\mathbf{u}$ & Vector CG2 \\
		Pressure space & $p$ & Scalar CG1 \\
		Adjoint formulation & -- & Frozen turbulence assumption \\
		Volume constraint & $V_{max}$ & [TBD]\% \\
		\hline
	\end{tabularx}
\end{table}

Resolving the shear layers on the inner radius of the extracted U-bend requires dense near-wall meshing. The results of the grid independence evaluation ensuring adequate spatial resolution are listed in Table~\ref{tab:app_ubend_mesh}.

\begin{table}[htbp]
	\centering
	\footnotesize
	\caption{Grid convergence study for the extracted U-Bend geometry (Ansys Fluent).}
	\label{tab:app_ubend_mesh}
	\renewcommand{\arraystretch}{1.2}
	\begin{tabularx}{\textwidth}{l c c c X}
		\hline
		\textbf{Model} & \textbf{Coarse Elements} & \textbf{Medium Elements} & \textbf{Fine Elements} & \textbf{Relative $\Delta p$ Change (Med-Fine)} \\
		\hline
		SA & [TBD] & [TBD] & [TBD] & [TBD]\% \\
		SST $k$-$\omega$ & [TBD] & [TBD] & [TBD] & [TBD]\% \\
		\hline
	\end{tabularx}
\end{table}

% ---------------------------------------------------------
% GLOBAL MESH CONVERGENCE FIGURES
% ---------------------------------------------------------
\section{Global Spatial Convergence Trends}
As indicated by the quantitative tables provided in the previous sections, the macroscopic pressure drop across all benchmark geometries exhibits clear asymptotic stabilization as spatial resolution increases. Visual confirmation of this global grid independence is aggregated in Figure~\ref{fig:app_opt_mesh_study}.

\begin{figure}[htpb]
	\centering
	\begin{subfigure}[b]{0.48\textwidth}
		\centering
		\meshplaceholder{Diffuser $\Delta p$ vs. Element Count}
		\caption{Diffuser}
		\label{fig:app_mesh_diffuser}
	\end{subfigure}
	\hfill
	\begin{subfigure}[b]{0.48\textwidth}
		\centering
		\meshplaceholder{Double-pipe $\Delta p$ vs. Element Count}
		\caption{Double pipe}
		\label{fig:app_mesh_doublepipe}
	\end{subfigure}
	
	\vspace{0.5cm}
	
	\begin{subfigure}[b]{0.48\textwidth}
		\centering
		\meshplaceholder{Curved benchmark $\Delta p$ vs. Element Count}
		\caption{Curved benchmarks}
		\label{fig:app_mesh_curved}
	\end{subfigure}
	\hfill
	\begin{subfigure}[b]{0.48\textwidth}
		\centering
		\meshplaceholder{Local velocity profile convergence at sensitive section}
		\caption{Profile-based convergence check}
		\label{fig:app_mesh_profiles}
	\end{subfigure}
	\caption{Asymptotic convergence of the macroscopic pressure drop metric and localized velocity profiles across successive mesh refinement levels for all optimized geometries.}
	\label{fig:app_opt_mesh_study}
\end{figure}

Across all evaluated geometries, the relative discrepancy between the Medium and Fine meshes ultimately falls below the strict numerical threshold of [Insert \%, e.g., 1.5\%]. This demonstrates that the governing equations are sufficiently resolved, ruling out spatial discretization error as a dominant factor. Consequently, the spatial configurations corresponding to the \textbf{Medium} mesh profiles were selected as the reliable standard for all comparative CFD validation campaigns executed in Chapter~\ref{ch:benchmark_results}.